\documentclass[11pt]{article}
\usepackage[margin=1in]{geometry}
\usepackage{amsmath,amssymb,amsfonts,amsthm}
\usepackage{bm}
\usepackage{hyperref}
\hypersetup{colorlinks=true, linkcolor=blue, urlcolor=blue, citecolor=blue}

\title{Glyph-Driven Codon Contrast:\\
A Falsifiable Bridge Between Symbolic Multiplicity and Biophysical Feature Space}
\author{Multiplicity Theory Working Notes}
\date{\today}

\begin{document}
\maketitle

\begin{abstract}
We describe a pilot framework that links prime-indexed glyph multiplicity to biophysical contrast features over coding sequences. On the biological side, we adopt the Codon-Contrast system: a 64-dimensional codon delta vector, a 64-dimensional Walsh--Hadamard biophysical contrast transform, and four flow metrics (transition ratio, GC bias per move, wobble-transition fraction, and flow entropy).\cite{codoncontrast} On the symbolic side, we use an 8-bit selector encoding glyph operations in a simplified commuting fragment (duplication, enclosure, base identity, medical-context flag). This selector is mapped through a resonance layer into 96 equivalence classes with exactly 32 classes of size 4 and 64 of size 2, enforcing a prime-like degeneracy structure. We provide a self-contained Python implementation and a permutation-based validation protocol over codon-aligned ortholog sets (e.g.\ TRPA1), designed to falsify or vindicate any claimed predictive link between glyph multiplicity and directions in biophysical contrast space.\cite{glyphengine}
\end{abstract}

\section{Executive Summary}

The goal is to test a sharp, falsifiable conjecture:
\emph{Do multiplicity features derived from ancient glyphs (in a prime-indexed operator sense) correspond to coherent directions in a low-dimensional, biophysically grounded feature space over genes targeted by the associated remedies?}\cite{glyphengine}

On the biology side, we fix the target category $\mathcal{B}$ as the Codon-Contrast feature space:
\begin{itemize}
  \item A 64-dimensional integer delta vector $\Delta h \in \mathbb{Z}^{64}$ capturing net codon count changes under a fixed lexicographic ordering of $\{A,C,G,T\}^3$;
  \item A 64-dimensional real vector of Walsh--Hadamard biophysical contrasts $w \in \mathbb{R}^{64}$ obtained by a length-64 FWHT after reindexing codons according to a 6-bit GC/PUR signature per position (GC1, PUR1, GC2, PUR2, GC3, PUR3), L1-normalized to yield scale-free contrasts;\cite{codoncontrast}
  \item Four scalar flow metrics summarizing substitution patterns: transition ratio, GC bias per move, wobble-transition fraction, and flow entropy.\cite{codoncontrast}
\end{itemize}

On the glyph side, we restrict attention to a commuting fragment of the glyph operator algebra (duplication and enclosure) and encode each glyph instance as an 8-bit selector $s\in\{0,1\}^8$ summarizing duplication, enclosure, stroke (here fixed), base glyph identity, and medical-context status. This selector is mapped via a deterministic resonance mapping to a class label $c(s)\in\{0,\dots,95\}$, with preimage counts enforcing a $32/64$ split between size-4 and size-2 classes.\cite{glyphengine}

The combined feature representation for a pair of codon-aligned sequences $(S_1,S_2)$ and a glyph selector $s$ is a hybrid vector
\[
  h_{\text{hyb}}(S_1,S_2,s) = \bigl(w(S_1,S_2),\, e(c(s))\bigr) \in \mathbb{R}^{160},
\]
where $w\in\mathbb{R}^{64}$ is the Walsh contrast vector and $e(c(s))\in\{0,1\}^{96}$ is a one-hot embedding of the resonance class. We then perform purely structural tests: first, a permutation-based assessment of non-random structure in Codon-Contrast features over ortholog sets; second, if structure exists, correlation tests between glyph-derived resonance classes and specific Codon-Contrast directions.

\section{Mathematical Overview}

\subsection{Biological Feature Space via Codon-Contrast}

Let $\Sigma = \{A,C,G,T\}$ and $\mathcal{C} = \Sigma^3$ denote the 64 codons, indexed lexicographically:
\begin{equation}
 \mathcal{C} = \{c_0, c_1,\dots, c_{63}\}.
\end{equation}
For a pair of codon-aligned coding sequences $S_1,S_2 \in \mathcal{C}^{L}$ (length $L$ codons), we define codon counts
\begin{equation}
  n_k^{(1)} = \#\{i : S_1[i] = c_k\},\qquad
  n_k^{(2)} = \#\{i : S_2[i] = c_k\}
\end{equation}
and the codon delta vector
\begin{equation}
  \Delta h_k = n_k^{(2)} - n_k^{(1)},\qquad k = 0,\dots,63.
\end{equation}
By construction, $\sum_{k=0}^{63} \Delta h_k = 0$, and $\Delta h$ encodes the exact net codon-level substitutions under a fixed frame and strand.\cite{codoncontrast}

To introduce biophysical structure, we map each codon $c\in\mathcal{C}$ to a 6-bit signature capturing GC and purine/pyrimidine properties at each position:
\begin{equation}
 \phi(c) = (\mathrm{GC}_1,\mathrm{PUR}_1,\mathrm{GC}_2,\mathrm{PUR}_2,\mathrm{GC}_3,\mathrm{PUR}_3) \in \{0,1\}^6,
\end{equation}
where, e.g., $\mathrm{GC}_i = 1$ if the $i$-th nucleotide is G or C, and $\mathrm{PUR}_i = 1$ if it is A or G.\cite{codoncontrast} This defines a permutation $\pi$ of $\{0,\dots,63\}$ such that indices are ordered by $\phi$.

Let $P$ be the $64\times 64$ permutation matrix implementing this reindexing, and let $H_{64}$ be the standard $64\times 64$ Hadamard matrix, recursively defined by
\begin{equation}
H_1 = \begin{bmatrix}1\end{bmatrix}, \quad
H_{2n} = \begin{bmatrix}H_n & H_n\\ H_n & -H_n\end{bmatrix}.
\end{equation}
The unnormalized Walsh--Hadamard biophysical contrasts are
\begin{equation}
  \tilde{w} = H_{64}\, P\, \Delta h \in \mathbb{R}^{64}.
\end{equation}
Let $\|\tilde{w}\|_1 = \sum_{k}|\tilde{w}_k|$. When $\|\tilde{w}\|_1>0$, we define the normalized contrast vector
\begin{equation}
  w = \frac{\tilde{w}}{\|\tilde{w}\|_1} \in \mathbb{R}^{64}.
\end{equation}
Each coordinate $w_k$ can be interpreted as a contrast along a particular combination of GC/PUR properties and their interactions, and the L1 normalization makes $w$ scale-free.\cite{codoncontrast}

In addition, four flow metrics summarize aggregate substitution structure over all positions:
\begin{itemize}
 \item Transition ratio:
       \[
       r_{\text{trans}} = \frac{\#\text{transitions }(A\leftrightarrow G, C\leftrightarrow T)}{\#\text{all substitutions}};
       \]
 \item Wobble-transition fraction:
       \[
       f_{\text{wobble}} = \frac{\#\text{transitions at third codon position}}{\#\text{all transitions}};
       \]
 \item GC bias per move:
       \[
       b_{\text{GC}} = \frac{\sum_{\text{substitutions}} (\mathrm{GC}(n_2) - \mathrm{GC}(n_1))}{\#\text{all substitutions}};
       \]
 \item Flow entropy:
       \[
       H_{\text{flow}} = -\sum_{(n_1\to n_2)} p_{(n_1\to n_2)} \log_2 p_{(n_1\to n_2)},
       \]
       where $p_{(n_1\to n_2)}$ is the empirical frequency of each nucleotide-level move.\cite{codoncontrast}
\end{itemize}

Thus for each pair $(S_1,S_2)$ we obtain a structured biological feature triple:
\begin{equation}
 (\Delta h, w, F_{\text{flow}}) \in \mathbb{Z}^{64} \times \mathbb{R}^{64} \times \mathbb{R}^4.
\end{equation}

\subsection{Glyph Selectors and Resonance Classes}

On the glyph side, we consider a commuting fragment of a prime-generated operator monoid acting on base glyphs. In the present pilot, we restrict to three observable operations:
\begin{itemize}
  \item Duplication $p_1$ (visual duplication of a glyph, interpreted as a dimerization-like operation);
  \item Enclosure $p_2$ (enclosing a glyph by a determinative or cartouche);
  \item A medical-context flag capturing whether the glyph appears in a medicinal context in the corpus.\cite{glyphengine}
\end{itemize}
Base glyphs are assigned discrete identities (e.g.\ garlic, willow, opium, etc.). For the 6-glyph pilot, we encode each glyph instance as an 8-bit selector
\begin{equation}
 s \in \{0,1\}^8,
\end{equation}
with bit assignments chosen as:
\begin{itemize}
  \item $b_0$: duplication present (1 if any $p_1$ applied);
  \item $b_1$: enclosure present (1 if any $p_2$ applied);
  \item $b_2$: stroke modification (fixed 0 in this fragment);
  \item $b_3$: medical-context flag (1 if glyph appears in a documented medical context);
  \item $b_4,b_5,b_6,b_7$: 4-bit base glyph identity code (0--15).
\end{itemize}
Thus each glyph instance (e.g.\ single garlic, duplicated garlic, etc.) is deterministically mapped to an 8-bit selector value.

We then define a resonance mapping
\begin{equation}
  \rho : \{0,1\}^8 \to \{0,\dots,95\} \times \{2,4\} \times \{0,1\},
\end{equation}
such that for each selector $s$ we obtain:
\begin{itemize}
  \item a class label $c(s)\in\{0,\dots,95\}$;
  \item a degeneracy $d(s)\in\{2,4\}$ (number of selectors in its equivalence class);
  \item a Boolean flag $\mathrm{size4}(s)$ indicating whether $d(s)=4$.
\end{itemize}
The mapping is constructed to satisfy:
\begin{equation}
 \#\{c : \#\rho^{-1}(c) = 4\} = 32,\quad
 \#\{c : \#\rho^{-1}(c) = 2\} = 64,
\end{equation}
and
\begin{equation}
 \sum_{c} \#\rho^{-1}(c) = 256.
\end{equation}
That is, the 256 possible selectors partition into 96 resonance classes, with an exact $32/64$ split in class degeneracy.\cite{glyphengine}

In code, this is implemented via a deterministic packing of certain bits (treated as ``free'' indices) and a compressed parameter $e\in\{-1,0,1\}$ derived from a specific pair of bits, but the key mathematical property is the enforced class-cardinality structure, checked at initialization.

For each selector $s$, we also define a one-hot embedding
\begin{equation}
 e(c(s)) \in \{0,1\}^{96},\quad e(c(s))_j =
 \begin{cases}
 1, & j=c(s),\\
 0, & \text{otherwise},
 \end{cases}
\end{equation}
and a normalized resonance index
\begin{equation}
 r_{\text{idx}}(s) = \frac{c(s)}{95} + \mathbf{1}\{d(s)=4\}\cdot \frac{1}{2},
\end{equation}
which serves as a combinatorial scalar coordinate (not yet claimed to be a biophysical invariant).\cite{glyphengine}

\subsection{Hybrid Feature Vector and Validation Strategy}

Given a codon-aligned sequence pair $(S_1,S_2)$ and an associated glyph selector $s$, we define the hybrid vector
\begin{equation}
 h_{\text{hyb}}(S_1,S_2,s) = \bigl(w(S_1,S_2),\, e(c(s))\bigr) \in \mathbb{R}^{160},
\end{equation}
which lives in a direct sum of the 64-dimensional contrast space and the 96-dimensional resonance class simplex. This provides a structured joint representation of biological substitution patterns and symbolic multiplicity class.

Our validation strategy is deliberately staged:
\begin{enumerate}
  \item \textbf{Structural test in Codon-Contrast space:} For a given gene family (e.g.\ TRPA1 orthologs), compute $(\Delta h, w, F_{\text{flow}})$ for all sequence pairs and evaluate whether the variance of each contrast coordinate and flow metric exceeds that obtained under a null model where codons are randomly permuted within each sequence. This is assessed by a permutation test yielding empirical p-values for each feature coordinate.\cite{codoncontrast}
  \item \textbf{Glyph--resonance overlay (only if Step 1 passes):} If the biophysical contrast space shows non-random structure, map the relevant glyphs (e.g.\ garlic, willow, opium glyphs) to selectors $s$ and classes $c(s)$, and then test whether resonance classes or indices correlate with specific contrast coordinates or flow metrics more strongly than codon-only baselines. Failure to find such correlations falsifies the strong glyph--biophysics functor hypothesis for the chosen operation fragment.
\end{enumerate}

\section{Core Python Implementation}

For concreteness, we include a compressed version of the key Python structures. The actual file is designed as a standalone script.

\subsection{Codon-Contrast and Flow Metrics}

\begin{verbatim}
NUCLEOTIDES = ['A','C','G','T']
CODONS = [a+b+c for a in NUCLEOTIDES for b in NUCLEOTIDES for c in NUCLEOTIDES]
CODON_INDEX = {c:i for i,c in enumerate(CODONS)}

def codon_delta(seq1, seq2):
    delta = np.zeros(64, dtype=int)
    for i in range(0, len(seq1), 3):
        c1 = seq1[i:i+3]; c2 = seq2[i:i+3]
        if c1 in CODON_INDEX and c2 in CODON_INDEX:
            delta[CODON_INDEX[c2]] += 1
            delta[CODON_INDEX[c1]] -= 1
    return delta

def compute_flow_metrics(seq1, seq2):
    total = transitions = gc_drift = wobble_transitions = 0
    move_counts = Counter()
    for i in range(0, len(seq1), 3):
        c1 = seq1[i:i+3]; c2 = seq2[i:i+3]
        for pos,(n1,n2) in enumerate(zip(c1,c2)):
            if n1 == n2: continue
            total += 1; move_counts[(n1,n2)] += 1
            if (n1,n2) in {('A','G'),('G','A'),('C','T'),('T','C')}:
                transitions += 1
                if pos == 2: wobble_transitions += 1
            gc_drift += (1 if n2 in 'GC' else 0) - (1 if n1 in 'GC' else 0)
    if total == 0:
        return dict(transition_ratio=0.0, wobble_fraction=0.0,
                    gc_bias=0.0, flow_entropy=0.0)
    tr = transitions / total
    wf = wobble_transitions / transitions if transitions else 0.0
    gc_bias = gc_drift / total
    H = 0.0
    for count in move_counts.values():
        p = count / total; H -= p * math.log2(p)
    return dict(transition_ratio=tr, wobble_fraction=wf,
                gc_bias=gc_bias, flow_entropy=H)

# H64 built recursively as a 64x64 Hadamard matrix
H64 = hadamard(64)

def codon_contrast(seq1, seq2):
    delta = codon_delta(seq1, seq2)
    walsh = H64 @ delta
    flow = compute_flow_metrics(seq1, seq2)
    return dict(delta_h=delta, walsh=walsh, flow=flow)
\end{verbatim}

\subsection{Resonance Mapping and Pilot Glyph Selectors}

\begin{verbatim}
class ResonanceMapper:
    def __init__(self):
        mapping = []
        for i in range(256):
            sel = f"{i:08b}"
            bits = [int(b) for b in sel]
            b0,b1,b2,b3,b4,b5,b6,b7 = bits
            free = (b1<<4)|(b2<<3)|(b3<<2)|(b6<<1)|b7
            if   (b4,b5) == (0,1): e = -1
            elif (b4,b5) == (1,0): e = 1
            else:                  e = 0
            e_offset = e + 1
            class_id = free * 3 + e_offset
            degeneracy = 4 if (b4 == b5) else 2
            is_size4 = (b4 == b5)
            mapping.append(dict(selector=sel, class_id=class_id,
                                degeneracy=degeneracy,
                                is_size4=is_size4,
                                resonance_index=class_id/95.0 +
                                               (0.5 if is_size4 else 0.0)))
        self.mapping = mapping
        class_ids = [m["class_id"] for m in mapping]
        assert len(set(class_ids)) == 96
        counts = Counter(class_ids)
        size4 = [cid for cid,cnt in counts.items() if cnt==4]
        size2 = [cid for cid,cnt in counts.items() if cnt==2]
        assert len(size4) == 32 and len(size2) == 64

    def from_selector(self, selector_str):
        for item in self.mapping:
            if item["selector"] == selector_str: return item
        raise ValueError("Selector not found")
\end{verbatim}

For the pilot, base glyphs (garlic, willow, opium, castor, royal, stroke) are assigned IDs in $\{0,\dots,5\}$, with selectors constructed as:
\begin{verbatim}
base_ids = {"garlic":0, "willow":1, "opium":2,
            "castor":3, "royal":4, "stroke":5}

def make_selector(base_name, dup=0, enc=0, stroke=0):
    base_id = base_ids[base_name]
    med = 1 if base_name in ["garlic","willow","opium","castor"] else 0
    val = (base_id<<4) | (med<<3) | (stroke<<2) | (enc<<1) | dup
    return f"{val:08b}"

single_glyphs = {base: make_selector(base,0,0) for base in base_ids}
\end{verbatim}

\section{Permutation Test Protocol}

Let $\{S^{(j)}\}_{j=1}^N$ be a set of codon-aligned ortholog sequences (e.g.\ TRPA1 in $N$ species).\cite{codoncontrast} For each unordered pair $(i,j)$, compute $w^{(ij)}$ and $F_{\text{flow}}^{(ij)}$. For each contrast coordinate $k$ and each flow metric $m$, compute the observed variance over pairs:
\begin{align}
  \sigma^2_{k,\text{obs}} &= \mathrm{Var}_{i<j}\bigl(w^{(ij)}_k\bigr),\\
  \sigma^2_{m,\text{obs}} &= \mathrm{Var}_{i<j}\bigl(F^{(ij)}_{\text{flow},m}\bigr).
\end{align}

To generate a null distribution, randomly permute codons within each $S^{(j)}$, recompute the features over all pairs, and record $\sigma^{2,(b)}_{k,\text{perm}}$ and $\sigma^{2,(b)}_{m,\text{perm}}$ for $b=1,\dots,B$ permutations. Empirical p-values are then
\begin{align}
  p_k &= \frac{1}{B}\sum_{b=1}^B \mathbf{1}\{\sigma^{2,(b)}_{k,\text{perm}}\ge \sigma^2_{k,\text{obs}}\},\\
  p_m &= \frac{1}{B}\sum_{b=1}^B \mathbf{1}\{\sigma^{2,(b)}_{m,\text{perm}}\ge \sigma^2_{m,\text{obs}}\}.
\end{align}
If few coordinates have $p<0.05$, the Codon-Contrast space is effectively noise for this gene, and any glyph→biology mapping should be regarded as falsified for that target. If multiple coordinates are significant, one can further test whether resonance classes $c(s)$ or resonance indices $r_{\text{idx}}(s)$ for appropriate glyphs explain additional variance or correlation beyond codon-only baselines.

\section{Conclusion and Outlook}

The present construction yields a fully specified, falsifiable bridge between symbolic multiplicity and codon-level biophysical contrasts. On the biological side, everything rests on the validated Codon-Contrast framework, which is rigorously defined and performance-tested but explicitly advertised as a feature engine rather than a fitness model.\cite{codoncontrast} On the symbolic side, we have imposed a strict 256$\to$96 class structure on glyph selectors, with exact degeneracy counts, to prevent unconstrained pattern-matching.\cite{glyphengine}

The key next steps are empirical: (i) run the permutation test on a curated TRPA1 (or other target) ortholog alignment to verify the existence of non-random structure in contrast space; (ii) if such structure is present, test whether glyph-derived resonance classes for documented medicinal glyphs (e.g.\ garlic, willow) correspond to stable directions in Codon-Contrast space. Success would motivate deeper work on the functorial relationship between prime-indexed glyph operators and biochemical transformations; failure would provide a sharp boundary on how far symbolic multiplicity can be pushed into biophysics.

\nocite{*}
\bibliographystyle{plainnat}
\bibliography{references}

\end{document}